% =====================================================================
\chapter{Lua-скрипты: программируем пульт}
\label{ch:lua}
% =====================================================================
\chapterintro
  {что такое Lua-скрипты, какие они бывают и где лежат; полезные готовые скрипты;
   как написать свою программу для экрана пульта}
  {SD-карта пульта, компьютер с текстовым редактором}

\section{Пульт, который можно программировать}

\begin{simple}
\textbf{Lua} --- простой язык программирования. EdgeTX умеет запускать Lua-программы
(скрипты) с SD-карты. Так сделаны меню настройки ExpressLRS, настройка Betaflight прямо с пульта,
разные мастера и даже игры. А ещё ты можешь написать свой скрипт --- например, красивый экран
с батареей.
\end{simple}

\begin{figure}[H]
\centering
\begin{tikzpicture}[every node/.style={font=\sffamily\small}]
  \foreach \x/\t/\f/\r in {0/Утилиты/TOOLS/{\btn{SYS} → Tools},3.9/Экраны/TELEMETRY/{кнопка \btn{TELE}},7.8/Функции/FUNCTIONS/{спецфункция\\«Lua script»},11.7/Миксеры/MIXES/{\btn{MDL} → Custom\\Scripts}}{
    \node[draw=etx,fill=etx!6,rounded corners=3pt,text width=32mm,minimum height=24mm,align=center] at (\x,0) {};
    \node[font=\sffamily\small\bfseries,text=etxdark] at (\x,0.75) {\t};
    \node[font=\ttfamily\scriptsize] at (\x,0.25) {SCRIPTS/\f};
    \node[font=\sffamily\scriptsize,align=center] at (\x,-0.5) {\r};}
\end{tikzpicture}
\caption{Четыре вида скриптов, их папки на SD-карте и как они запускаются}
\end{figure}

\begin{center}\small
\renewcommand{\arraystretch}{1.2}
\begin{tabularx}{\linewidth}{@{}L{26mm}Y@{}}
\toprule
Утилиты & запускаются из \mpath{SYS\sep Tools}, занимают весь экран, пока не выйдешь \\
Экраны телеметрии & показываются по кнопке \btn{TELE} \\
Функции & работают, пока включён переключатель в специальной функции \\
Миксеры & считают своё значение, которое можно использовать в миксерах \\
\bottomrule
\end{tabularx}
\end{center}

\section{Готовые полезные скрипты}

\begin{description}[font=\sffamily\bfseries\color{etx}]
  \item[ExpressLRS] настройка модуля и приёмника (глава~\ref{ch:rf}).
  \item[Betaflight TX Lua Scripts] настройка PID, рейтов, видеопередатчика квадрокоптера
        прямо с пульта. Скачиваются с GitHub проекта Betaflight, копируются в папку
        \code{SCRIPTS}, запускаются из \menu{Tools}.
  \item[Мастера моделей] входят в пакет SD-карты EdgeTX.
\end{description}

\section{Экран пульта --- это клетчатый лист}

Экран TX12 и Boxer --- 128 точек в ширину и 64 в высоту. Координата $x$ растёт вправо,
$y$ --- \textbf{вниз} (не вверх, как в школе!). Точка $(0,0)$ --- левый верхний угол.

\begin{figure}[H]
\centering
\begin{tikzpicture}[x=0.78mm,y=0.78mm]
  \fill[lcdbg] (0,0) rectangle (128,-64);
  \foreach \x in {0,16,...,128}{\draw[lcdfg!25] (\x,0) -- (\x,-64);}
  \foreach \y in {0,-16,...,-64}{\draw[lcdfg!25] (0,\y) -- (128,\y);}
  \draw[lcdfg,line width=0.8pt] (0,0) rectangle (128,-64);
  \foreach \x in {0,32,64,96,128}{\node[font=\sffamily\scriptsize,above] at (\x,0) {\x};}
  \foreach \y in {0,16,32,48,64}{\node[font=\sffamily\scriptsize,left] at (0,-\y) {\y};}
  \draw[->,accent,thick] (0,6) -- (20,6) node[right,font=\sffamily\scriptsize]{$x$};
  \draw[->,accent,thick] (-9,0) -- (-9,-16) node[below,font=\sffamily\scriptsize]{$y$};
  % пример
  \fill[lcdfg] (0,0) rectangle (60,-8);
  \node[anchor=north west,font=\ttfamily\scriptsize,text=lcdbg,inner sep=1pt] at (1,-0.5) {Quad-5in};
  \node[anchor=north west,font=\ttfamily\bfseries\LARGE,text=lcdfg,inner sep=0pt] at (0,-12) {3.95};
  \draw[lcdfg,line width=0.6pt] (0,-32) rectangle (80,-42);
  \fill[lcdfg] (1,-33) rectangle (60,-41);
  \node[anchor=north west,font=\ttfamily\scriptsize,text=lcdfg,inner sep=0pt] at (86,-12) {LQ 100};
  \fill[accent] (0,0) circle (1.3);
  \node[font=\sffamily\scriptsize,text=accent,anchor=west] at (131,-1) {(0,0) --- начало};
  \node[font=\sffamily\scriptsize,anchor=west] at (131,-14) {\code{drawNumber(0, 12, …)}};
  \node[font=\sffamily\scriptsize,anchor=west] at (131,-26) {\code{drawText(86, 12, "LQ")}};
  \node[font=\sffamily\scriptsize,anchor=west] at (131,-38) {\code{drawRectangle(0, 32, 80, 10)}};
\end{tikzpicture}
\caption{Координаты экрана и что рисуют команды}
\end{figure}

\section{Как устроен скрипт}

Скрипт --- это файл \code{.lua}, который в конце \emph{возвращает таблицу} с функциями:
\begin{itemize}
  \item \code{init()} --- выполняется один раз при запуске;
  \item \code{run(...)} --- выполняется снова и снова, много раз в секунду: тут рисуем экран.
\end{itemize}

Самые нужные команды:
\begin{center}\small
\renewcommand{\arraystretch}{1.2}
\begin{tabularx}{\linewidth}{@{}L{58mm}Y@{}}
\toprule
\code{getValue("RxBt")} & узнать значение датчика или источника \\
\code{lcd.clear()} & очистить экран \\
\code{lcd.drawText(x, y, "текст", флаги)} & написать текст; флаги: \code{INVERS} (инверсия),
  \code{BLINK} (мигать), \code{SMLSIZE}, \code{MIDSIZE}, \code{DBLSIZE} (размер) \\
\code{lcd.drawNumber(x, y, число, флаги)} & написать число; \code{PREC1}/\code{PREC2} ---
  1 или 2 знака после запятой \\
\code{lcd.drawRectangle(x, y, w, h)} & рамка \\
\code{lcd.drawFilledRectangle(x, y, w, h)} & закрашенный прямоугольник \\
\code{playNumber(n, единица)} & сказать число голосом \\
\code{getTime()} & текущее время в «тиках» по 10\,мс \\
\bottomrule
\end{tabularx}
\end{center}

\section{Пример 1: «Привет, EdgeTX!»}

\luafile{hello.lua}

Сохрани как \code{SCRIPTS/TOOLS/hello.lua}, открой \mpath{SYS\sep Tools} и запусти
\menu{hello}. Нажимай \btn{ENTER} --- счётчик растёт; \btn{RTN} --- выход.

\begin{center}
\lcdscreen{Hello, EdgeTX!}{}{\lblank\lr{Model: Quad-5in}{}\lr{ENTER pressed:}{3\hspace{12mm}}\lblank\lr{Tx batt: 7.9V}{}}
\end{center}

\section{Пример 2: экран батареи}

\luafile{batt.lua}

Файл кладётся в \code{SCRIPTS/TELEMETRY/batt.lua} (имя --- не длиннее 6 букв). Затем
\mpath{MDL\sep Telemetry\sep Screen 1} = \menu{Script} → \menu{batt}. Нажми \btn{TELE}:

\begin{figure}[H]
\centering
\begin{tikzpicture}[x=0.55mm,y=0.55mm]
  \fill[black!80,rounded corners=3mm] (-7,7) rectangle (135,-71);
  \fill[lcdbg] (0,0) rectangle (128,-64);
  \fill[lcdfg] (0,0) rectangle (42,-8);
  \node[anchor=north west,font=\ttfamily\scriptsize,text=lcdbg,inner sep=1pt] at (1,-0.5) {Quad-5in};
  \node[anchor=north west,font=\ttfamily\bfseries\Large,text=lcdfg,inner sep=0pt] at (0,-12) {3.95};
  \node[anchor=north west,font=\ttfamily\tiny,text=lcdfg,inner sep=0pt] at (38,-19) {V/cell};
  \draw[lcdfg,line width=0.5pt] (0,-32) rectangle (80,-42);
  \fill[lcdfg] (1,-33) rectangle (56,-41);
  \node[anchor=north west,font=\ttfamily\tiny,text=lcdfg,inner sep=0pt] at (86,-12) {LQ  100};
  \node[anchor=north west,font=\ttfamily\tiny,text=lcdfg,inner sep=0pt] at (86,-24) {RSSI -54};
  \node[anchor=north west,font=\ttfamily\tiny,text=lcdfg,inner sep=0pt] at (0,-48) {Pack 15.8};
  \node[anchor=north west,font=\ttfamily\tiny,text=lcdfg,inner sep=0pt] at (86,-48) {02:37};
\end{tikzpicture}
\caption{Так выглядит экран, нарисованный скриптом \code{batt.lua}}
\end{figure}

\section{Пример 3: скрипт-миксер}

\luafile{thrlim.lua}

Скрипт берёт стик газа и ограничивает его. \mpath{MDL\sep Custom Scripts} → \menu{thrlim},
укажи источник \sw{Thr} и предел. В миксере газа используй источник \sw{Lim}. Удобно для
новичков: «не больше 60\,\% газа».

\section{Пример 4: говорящий скрипт}

\luafile{lqcall.lua}

Кладётся в \code{SCRIPTS/FUNCTIONS/lqcall.lua}. В специальных функциях: \menu{Switch} =
\sw{SA↓}, \menu{Lua script} = \menu{lqcall}. В арме пульт каждые 10 секунд называет
качество связи.

\section{Если что-то не работает}

\begin{itemize}
  \item \code{Script syntax error} --- ошибка в тексте программы (забытая скобка, \code{end}).
  \item Скрипта нет в списке --- проверь папку и длину имени файла.
  \item Проверяй скрипты на компьютере: в EdgeTX Companion есть симулятор пульта, который
        запускает Lua и показывает сообщения \code{print()}.
\end{itemize}

\begin{tryit}
\begin{enumerate}
  \item Запусти \code{hello.lua}. Поменяй текст приветствия на своё имя.
  \item В \code{batt.lua} поменяй \code{CELLS} под свою батарею и выведи экран по \btn{TELE}.
  \item Переделай \code{lqcall.lua}, чтобы он говорил напряжение \sw{RxBt}.
\end{enumerate}
\end{tryit}

\begin{summary}
\begin{itemize}
  \item Скрипты лежат в папке \code{SCRIPTS} на SD-карте: \code{TOOLS}, \code{TELEMETRY},
        \code{FUNCTIONS}, \code{MIXES}.
  \item Экран 128×64: $x$ вправо, $y$ вниз, $(0,0)$ --- левый верхний угол.
  \item \code{getValue()} --- прочитать датчик, \code{lcd.draw…()} --- нарисовать.
\end{itemize}
\end{summary}
